\documentclass[12pt,a4paper]{article}
% Preambule commun aux comptes rendus CR_*.tex.

\usepackage{array}
\newcolumntype{C}[1]{>{\centering\arraybackslash}m{#1}}
\usepackage{multirow}
\usepackage{geometry}
\geometry{a4paper, margin=1in}
\usepackage{graphicx}
\usepackage{float}
\usepackage{tikz}
\usetikzlibrary{arrows.meta,positioning}
\usepackage{pgfgantt}
\usepackage{pdflscape}
\usepackage{enumitem}
\usepackage{booktabs}
\usepackage{longtable}
\usepackage{ragged2e}
\usepackage{tabularx}
\usepackage{xparse}
\usepackage{ltablex}
\usepackage{xltabular}
\usepackage{subcaption}

\newcolumntype{Y}{>{\RaggedRight\arraybackslash}X}

\newcommand{\StyledTableSetup}{%
  \footnotesize
  \renewcommand{\arraystretch}{1.2}%
  \setlength{\tabcolsep}{4pt}%
}


% ============================================================
% IHEDN COVER TEMPLATE (XeLaTeX)
% ============================================================
% Compilation (from project root):
%   latexmk -pdf main.tex
%   LATEXMK_SHELL_ESCAPE=0 latexmk -pdf main.tex
%
% Optional environment controls from latexmkrc:
%   LATEXMK_ENGINE=xelatex        (default/enforced)
%   LATEXMK_SHELL_ESCAPE=1|0      (default 1, needed for SVG conversion)
%
% Required package for this template:
%   \usepackage{ihedn-cover}
%
% Note: ihedn-cover already loads needed internal packages
% (fontspec, graphicx, tikz, ragged2e, optional svg).
% ============================================================

% ----------------------------------------------------------------
% Package options (documented)
% ----------------------------------------------------------------
% Geometry scope:
%   apply-geometry            -> force A4 + zero margins while cover is built
%   no-apply-geometry         -> do not alter host geometry (default)
%
% Typesetting freeze:
%   global-typesetting        -> apply freeze globally (intrusive)
%   no-global-typesetting     -> no global freeze (default)
%
% Small-caps handling:
%   local-smallcaps-fix       -> \textsc -> \MakeUppercase only inside cover (default)
%   no-local-smallcaps-fix    -> disable local fix
%   global-smallcaps-fix      -> global override (intrusive)
%   no-global-smallcaps-fix   -> disable global fix (default)
%
% Main font handling:
%   global-mainfont           -> set Times New Roman globally (default)
%   local-mainfont            -> set Times only during cover rendering
%   no-mainfont-override      -> keep host document main font
%
% SVG support:
%   svg-support               -> enable SVG logos (default, needs shell-escape)
%   no-svg-support            -> disable SVG path handling (pdf/png only)
%
% Debug overlay:
%   debug                     -> show mm grid + bounding boxes on cover
%   no-debug                  -> hide debug overlay (default)
% ----------------------------------------------------------------
\usepackage[
  no-apply-geometry,
  no-global-typesetting,
  global-smallcaps-fix,
  local-smallcaps-fix,
  global-mainfont,
  svg-support,
  no-debug
]{ihedn-cover}

% ----------------------------------------------------------------
% Bibliography stack (BibLaTeX/Biber, style from subtree vendor/)
% ----------------------------------------------------------------
\usepackage{polyglossia}
\setdefaultlanguage{french}
\makeatletter
\g@addto@macro\captionsfrench{%
  \renewcommand{\tablename}{Tableau}%
  \renewcommand{\figurename}{Figure}%
}
\makeatother
\usepackage{csquotes}
\usepackage{xurl}
\usepackage{hyperref}
\makeatletter
\let\ihedn@orig@input@path\input@path
\def\input@path{{vendor/biblatex-sciences-po-fr/style/}}
\makeatother
\usepackage[
  backend=biber,
  style=sciencespo,
  hyperref=true,
  autolang=other
]{biblatex}
\AtBeginBibliography{\sloppy\setlength{\emergencystretch}{3em}}
\makeatletter
\let\input@path\ihedn@orig@input@path
\makeatother
\addbibresource{bibliographies/bib/rapport.bib}

% ----------------------------------------------------------------
% tcolorbox
% ----------------------------------------------------------------
\usepackage[most]{tcolorbox}
\tcbuselibrary{breakable}

% Example: use Arial for the full report + cover.
% Uncomment and switch package option to no-mainfont-override.
% \setmainfont{Arial}[
%   UprightFont    = *,
%   BoldFont       = * Bold,
%   ItalicFont     = * Italic,
%   BoldItalicFont = * Bold Italic
% ]

%==================================================
% Liste des propositions
%==================================================

\makeatletter

\newcommand{\listpropositionname}{Liste des propositions}

\newcommand{\listofpropositions}{
  \section*{\listpropositionname}
  \addcontentsline{toc}{section}{\listpropositionname}
  \@starttoc{lop}
}

\newcommand*\l@proposition{\@dottedtocline{1}{1.5em}{2.3em}}

\makeatother

%==================================================
% Compteur des propositions
%==================================================

\newcounter{proposition}

%==================================================
% Style graphique des propositions
%==================================================

\newtcolorbox{propositionbox}[2][]{
  lower separated=false,
  colback=white!80!gray,
  colframe=white!20!black,
  fonttitle=\bfseries,
  colbacktitle=white!30!gray,
  coltitle=black,
  enhanced,
  sharp corners,
  attach boxed title to top left={
    xshift=0.5cm,
    yshift=-2mm
  },
  title=#2,
  #1
}

%==================================================
% Environnement proposition
%==================================================

\newenvironment{proposition}[1]
{
  \refstepcounter{proposition}

  \addcontentsline{lop}{proposition}{%
    \protect\numberline{\theproposition}#1%
  }

  \begin{propositionbox}
  {Proposition~\theproposition~: #1}
}
{
  \end{propositionbox}
}

% Renommage Liste des figures en Table des figures
\addto\captionsfrench{%
  \renewcommand{\listfigurename}{Table des figures}
}

\newcommand{\heure}[2]{{\NoAutoSpacing #1:#2}}

\DeclareFieldFormat[misc]{title}{#1}

% Activer ce package avec l'option none retire la césure.
%\usepackage[none]{hyphenat}

\usepackage{adjustbox}
\usepackage{pdfpages}

\makeatletter
\newcommand{\IhednSetPdfPageCount}[2]{%
  \begingroup
  \edef\ihedn@pdfpath{#2}%
  \ifXeTeX
    \global#1=\XeTeXpdfpagecount "\ihedn@pdfpath" %
  \else
    \ifLuaTeX
      \saveimageresource{\ihedn@pdfpath}%
      \global#1=\lastsavedimageresourcepages\relax
    \else
      \pdfximage{\ihedn@pdfpath}%
      \global#1=\pdflastximagepages\relax
    \fi
  \fi
  \endgroup
}
\makeatother

\newcount\IhednIncludedPdfPageCount
\newcommand{\IhednRenderPdfPagesInBoxes}[1]{%
  \IhednSetPdfPageCount{\IhednIncludedPdfPageCount}{#1}%
  \foreach \IhednIncludedPdfPage in {1,...,\the\IhednIncludedPdfPageCount}{%
    \begin{tcolorbox}[
      colback=gray!5,
      colframe=black,
      boxrule=0.5pt,
      arc=3mm
    ]
    \centering
    \includegraphics[page=\IhednIncludedPdfPage,width=0.95\textwidth]{#1}
    \end{tcolorbox}
    \ifnum\IhednIncludedPdfPage<\IhednIncludedPdfPageCount
      \clearpage
    \fi
  }%
}

\usepackage{tablefootnote}


\begin{document}

% ----------------------------------------------------------------
% Cover keys (all available keys are shown below)
% ----------------------------------------------------------------
% Header block (top-right), logo, title, report lines, subtitle, committee.
\PageDeGardeIHEDN{
    header_1={Union-IHEDN},
    header_2={Association des auditeurs IHEDN},
    header_3={région Paris / Île de France},
    date={le 19 décembre 2025},
    logo_path={Figures/logos_IHEDN/Logo_AR_16_PARIS_IDF.jpg},
    title={Crises majeures : quel rôle pour les citoyens engagés et les associations IHEDN ?},
    title_max_lines=2,
    title_fontsize={26.000pt},
    title_baselineskip={28.667pt},
    reportline_1={Rapport du Comité d'étude de l'Association régionale des auditeurs},
    reportline_2={de l'IHEDN région Paris / Île de France (AR 16)},
    subtitle={Compte rendu de la réunion du 11 décembre 2025},
    committee={Comité d'étude, d'octobre 2025 à juin 2026, composé de : Mme Valérie Arlé-Faivre, \textit{présidente} ; M. Aurélien Duvignac-Rosa, \textit{rapporteur} ; M. Pascal Bayot Duponchel ; M. Jean-Jacques Cleynen ; Mme Valérie Cowan ; M. Alain Fontan ; M. Pierre-Marie Giard ; Mme Muriel Joyeux ; M. Bruno Leray ; M. Yves-Henri Renhas ; M. Marc Roure ; Mme Isabelle de Segonzac ; M. Gérard Turck ; Mme Marie Danielle Vasquez-Duchêne.}
}

\section{Rappel des échanges antérieurs et recentrage des travaux}

Lors de la première réunion, les travaux du comité avaient principalement porté sur les enjeux de la guerre informationnelle, en particulier la lutte contre les infox\footnote{cf. \url{https://www.culture.fr/franceterme/terme/CULT754} [consulté le 19 déc. 2025].} et les phénomènes de désinformation.  
À l'issue des échanges intervenus lors de la réunion du 11 décembre 2025, il a été décidé de recentrer les travaux du comité sur le rôle effectif que peuvent jouer les citoyens engagés et les associations régionales de l'IHEDN dans la prévention, la gestion et la résilience face aux crises majeures.

Les objectifs désormais poursuivis sont les suivants :

\begin{itemize}
    \item Définir les modalités d'implication des citoyens et des auditeurs de l'IHEDN dans la gestion des crises (attentats, catastrophes naturelles, espionnage industriel, etc.) ;
    \item Renforcer la coordination entre les acteurs locaux (mairies, préfectures, forces de sécurité) et les associations IHEDN ;
    \item Identifier des leviers opérationnels permettant d'améliorer la préparation et la réponse aux crises, en s'appuyant sur les compétences et l'expertise des membres de l'Institut.
\end{itemize}

\section{Constats et enjeux majeurs}

\subsection{Typologie des crises et géographie des risques}

Les participants ont souligné la nécessité de prendre en compte une diversité de situations de crise, notamment :

\begin{itemize}
    \item Les crises informationnelles (infox, manipulation de l'opinion) ;
    \item les crises sécuritaires (attentats, menaces terroristes) ;
    \item les crises industrielles (espionnage, cyberattaques) ;
    \item les crises sanitaires et environnementales (pandémies, catastrophes naturelles).
\end{itemize}

Il a été rappelé que chaque territoire présente des spécificités propres, impliquant des réponses adaptées, tant au niveau national qu'au niveau local.

\subsection{Rôle des citoyens et des associations IHEDN}

Les retours d'expérience ont mis en évidence :

\begin{itemize}
    \item Une forte volonté de participation et d'implication citoyenne, souvent fondée sur un registre émotionnel réactif et peu structuré, tel qu'il s'est manifesté lors de crises antérieures. Cette mobilisation demeure toutefois limitée par un déficit de visibilité d'ensemble et par une connaissance insuffisante des dispositifs et plans existants ;
    \item Dans un contexte de crise appelé à s'intensifier\footnote{conformément aux éléments exposés par le général Hervé de Courrèges lors du dîner de rentrée de l'AR 16.}, il apparaît nécessaire d'organiser une mobilisation plus réactive des associations de l'IHEDN, dont les membres disposent de compétences diversifiées et reconnues ;
    \item L'absence d'un cadre clairement défini ne permet pas une mobilisation efficace et coordonnée des auditeurs en appui des autorités publiques.
\end{itemize}

\subsection{Insuffisances et points d'attention identifiés}

Plusieurs insuffisances ont été identifiées, notamment :

\begin{itemize}
    \item un manque de coordination entre les différents acteurs concernés (citoyens, associations, institutions) ;
    \item l'absence d'un socle commun de formation pour les membres de l'IHEDN susceptibles d'intervenir en situation de crise ;
    \item la disparition de la réserve IHEDN, dont il conviendrait d'analyser les causes et d'évaluer l'opportunité d'une réactivation sous une forme adaptée aux enjeux actuels.
\end{itemize}

\section{Propositions et axes d'action}

\subsection{Mobilisation et valorisation des compétences des auditeurs}

\textbf{Cartographie des ressources}

\renewcommand{\arraystretch}{1.3}
\begin{center}
\begin{tabular}{|C{4cm}|C{6cm}|C{4cm}|}
\hline
\textbf{Axe d'action} & \textbf{Proposition} & \textbf{Objectif} \\
\hline
Cartographie des compétences &
Élaborer et diffuser un questionnaire auprès des membres de l'IHEDN afin
d'identifier leurs compétences, leur disponibilité et leurs zones potentielles
d'intervention\footnotemark &
Disposer d'une vision consolidée des ressources mobilisables. \\
\hline
Base de données centralisée &
Constituer une base de données regroupant les compétences identifiées, mobilisable
par les autorités locales en cas de crise. &
Faciliter la mobilisation rapide et coordonnée des auditeurs. \\
\hline
Réserve IHEDN &
Étudier les conditions de mise en place d'une réserve opérationnelle IHEDN, en lien
avec les préfets et les mairies. &
Apporter un appui logistique, technique ou humain en situation de crise. \\
\hline
Analyse du dispositif antérieur &
Analyser les raisons ayant conduit à l'interruption de la précédente réserve IHEDN. &
Tirer les enseignements nécessaires avant toute réactivation. \\
\hline
Identification des volontaires &
Identifier les membres volontaires afin de garantir leur légitimité et leur
visibilité auprès des autorités comme de la population. &
Assurer la crédibilité et l'efficacité du dispositif. \\
\hline
\end{tabular}
\end{center}

\footnotetext{Proposition portée par M. Bruno Leray et rapportée lors de la réunion par M. Alain Fontan.}

\newpage

\subsection{Sensibilisation et formation}

L'organisation de la gestion de crise en France repose sur une structuration à plusieurs niveaux. Dans ce cadre, différents axes d'action peuvent être envisagés afin de favoriser les échanges autour des enjeux et problématiques liés à la gestion de crise.

\renewcommand{\arraystretch}{1.3}
\begin{center}
\begin{tabular}{|C{3.5cm}|C{5.5cm}|C{3.5cm}|C{3.5cm}|}
\hline
\textbf{Axe d'action} & \textbf{Proposition} & \textbf{Objectif} & \textbf{Intervenants} \\
\hline
\multirow{3}{*}{\parbox[c][\totalheight][c]{3.5cm}{\vspace{2.2cm} \centering Sessions d'information}} &
Présenter en conseils municpaux les dispositifs de plans communaux de sauvegarde (PCS) &
Renforcer la connaissance des dispositifs de gestion de crise au niveau local. &
Élus locaux \\
\cline{2-4}
&
Mettre en place des actions de sensibilisation à destination des conseils municipaux
des enfants, fondées sur une approche pédagogique adaptée à leur tranche d'âge. &
Développer une culture de résilience dès le plus jeune âge. &
Élus locaux, enseignants, intervenants spécialisés \\
\cline{2-4}
&
Organiser des soirées thématiques à destination des élus locaux (maires notamment) et des sessions d'information pilotées par les préfets à destination des élus
locaux, portant sur les dispositifs de gestion de crise. &
Renforcer la coordination et l'appropriation des dispositifs de crise au niveau local. &
Préfets \\
\hline
Partenariats académiques et médiatiques, appuyés par l'IHEDN &
Impliquer les trinômes académiques et les écoles de journalisme afin de renforcer
la culture de sécurité et la lutte contre les infox. &
Développer une culture partagée de vigilance et d'analyse de l'information. &
Auditeurs IHEDN, Universitaires, formateurs, professionnels des médias \\
\hline
Modules de formation &
Collaborer avec le CLEMI\footnotemark[3] et la plateforme VIGINUM\footnotemark[4]
pour élaborer des modules de formation relatifs à la gestion de crise et à la
détection des infox. &
Structurer une offre de formation adaptée aux enjeux contemporains. &
CLEMI, VIGINUM, experts en cybersécurité et information stratégique \\
\hline
Intégration aux parcours IHEDN &
Intégrer les modules de formation développés aux parcours de formation des
auditeurs de l'IHEDN. &
Assurer une montée en compétence progressive et homogène des auditeurs. &
Corps enseignant IHEDN, intervenants institutionnels \\
\hline
\end{tabular}
\end{center}

\footnotetext[3]{Centre pour l'éducation aux médias et à l'information (cf. \url{https://www.clemi.fr/} [consulté le 19 déc. 2025]).}
\footnotetext[4]{VIGINUM est le service technique et opérationnel chargé de la vigilance et de la protection contre les ingérences numériques étrangères (cf. \url{https://www.sgdsn.gouv.fr/notre-organisation/composantes/service-de-vigilance-et-protection-contre-les-ingerences-numeriques} [consulté le 19 déc. 2025]).}

\subsection{Renforcement de la coordination avec les autorités locales}

\paragraph{Référencement institutionnel}
~\\
Il a été proposé :

\begin{itemize}
    \item de formaliser des partenariats avec les préfectures, les mairies et les forces de sécurité afin que les associations IHEDN soient identifiées comme des acteurs ressources en cas de crise ;
    \item d'intégrer ces associations aux dispositifs locaux de gestion de crise (plans de secours, cellules de crise, etc.).
\end{itemize}

\paragraph{Exercices et simulations}
~\\
L'organisation d'exercices et de simulations associant les acteurs locaux a été évoquée, afin de tester les dispositifs existants et d'améliorer la réactivité collective.

\newpage

\section{Points de vigilance et défis à relever}

\begin{itemize}
    \item \textbf{Acceptabilité locale} : certaines réticences pourraient être exprimées par les élus locaux à des fins politique, l'intervention des associations IHEDN pouvant être perçue comme une ingérence ou une complexification des dispositifs existants ;
    \item \textbf{Garantie des compétences} : il apparaît indispensable de s'assurer que les membres mobilisés disposent des connaissances, outils et formations nécessaires ;
    \item \textbf{Disponibilité des auditeurs} : les contraintes professionnelles et personnelles des membres devront être prises en compte afin de garantir un engagement réaliste et durable ;
    \item \textbf{Communication} : une attention particulière devra être portée à la clarté, à la transparence et à la pédagogie de la communication à destination du public via divers moyens notamment médiatiques.
\end{itemize}

\section{Éléments à retenir à l'issue de la réunion}

\begin{center}
\renewcommand{\arraystretch}{1.3}
\begin{tabular}{|p{4cm}|p{10cm}|}
\hline
\textbf{Thématique} & \textbf{Éléments clefs} \\
\hline
Cadrage général &
Recentrage des travaux du comité sur le rôle des citoyens engagés et des associations régionales de l'IHEDN dans la prévention, la gestion et la résilience face aux crises majeures. \\
\hline
Constats partagés &
Existence de compétences significatives au sein des réseaux IHEDN, mais nécessité de coordination, de formation et de cadre d'intervention structuré. \\
\hline
Orientations proposées &
Cartographie des compétences des auditeurs, réflexion sur la réactivation d'une réserve IHEDN adaptée aux enjeux actuels, renforcement des partenariats avec les autorités locales. \\
\hline
Conditions de réussite &
Acceptabilité par les élus et autorités locales, garantie du niveau de compétence des membres mobilisés, prise en compte réaliste de la disponibilité des auditeurs. \\
\hline
Points de vigilance &
Nécessité d'une communication claire, pédagogique, simple et transparente afin d'assurer la compréhension et la légitimité des actions auprès de la population. \\
\hline
\end{tabular}
\end{center}

\newpage

\section{Calendrier des prochaines réunions}

\begin{table}[h!]
\centering
\renewcommand{\arraystretch}{1.3}
\begin{tabular}{|l|l|l|}
\hline
\textbf{Jour} & \textbf{Date} & \textbf{Salle} \\
\hline
Jeudi & 22 janvier 2026 & Salle Castex \\
Jeudi & 12 février 2026 & Salle Castex \\
Jeudi & 5 mars 2026 & Salle Castex \\
Jeudi & 19 mars 2026 & Salle Castex \\
Jeudi & 9 avril 2026 & Salle Castex \\
Jeudi & 23 avril 2026 & Salle Castex \\
Jeudi & 21 mai 2026 & Salle Aubry \\
Jeudi & 11 juin 2026 & Salle Castex \\
\hline
\end{tabular}
\end{table}

\end{document}